\begin{table*}[htp]
\resizebox{\textwidth}{!}{
\begin{tabular}{cll|lllll}
\toprule
& \multicolumn{2}{c|}{\textbf{Reiteration}} & \multicolumn{5}{c}{\textbf{Wedging}} \\
Model & Quality  & Style & Qual. 1 & Qual. 2 & Qual. 3 & Style & Hostility \\
\midrule
\anthropic & 3.975 (0.892) & 4.343 (0.659) & 0.364 (0.703) & 0.333 (0.711) & 0.515 (0.520) & 0.848 (0.261) & 0.848 (0.702)\\
\optonesevenfive & 3.814 (0.841) & 4.314 (0.557) & 0.121 (0.879) & 0.545 (0.608) & 0.273 (0.664) & 0.879 (0.257) & 0.348 (0.484)\\
\optsixsix & 3.426 (0.993) & 2.990 (1.297) & -0.061 (0.789) & -0.000 (0.804) & -0.152 (0.702) & 0.424 (0.494) & 0.242 (0.378)\\
\gptdavinci & 3.598 (0.860) & 4.113 (0.797) & 0.212 (0.608) & 0.485 (0.539)& 0.152 (0.744) & 0.606 (0.509) & 0.500 (0.762)\\
\instructdavinci & 4.221 (0.779) & 4.407 (0.498) & 0.273 (0.814) & 0.727 (0.467) & 0.212 (0.456) & 0.939 (0.192) & 0.485 (0.641)\\
\glm & 3.946 (0.781) & 1.270 (0.499) & 0.364 (0.758) & 0.364 (0.731) & 0.303 (0.731) & -0.576 (0.514) & 0.727 (0.664)\\
\bottomrule
\end{tabular}}
\caption{\textbf{Human evaluation for disinformation scenarios.}
Note: Qual. 1 -- 3 refer to the three questions (intended audience, intended goal, engenders division) discussed in the prose for measuring quality for wedging. Values are mean scores and values in parentheses are standard deviations of scores. Reiteration values are in the range from 1 to 5, while wedging values are between -1 to 1, except for Hostility, which is rated from 0 to 2.}
\label{tab:disinformation-human-evaluation}
\end{table*}